\section{Напредни примери}
Ови примери претпостављају познавање основних концепата TLA+ и приказују реалистичније системе
са вишеструким процесима, збиркама порука и нетривијалним особинама исправности.

\subsection{Outbox pattern}
\begin{primer}[Outbox pattern]
Outbox pattern је техника којом се у истој трансакцији уписује промена у базу и порука у \texttt{outbox} табелу.
Засебан процес затим чита \texttt{outbox} и шаље поруке ка брокеру, па се порука чисти тек након потврде.

\textbf{TLA+ код (издвојени део спецификације):}
\begin{lstlisting}
VARIABLES db, outbox, broker, sent, acked, allMsgs, nextMsgId

CreateOrder(o) ==
    LET m == [id |-> nextMsgId, order |-> o] IN
    /\ db[o] = "NEW"
    /\ db' = [db EXCEPT ![o] = "CREATED"]
    /\ outbox' = outbox \cup {m}
    /\ allMsgs' = allMsgs \cup {m}
    /\ nextMsgId' = nextMsgId + 1
    /\ UNCHANGED <<broker, sent, acked>>

RelaySend(m) ==
    /\ m \in outbox
    /\ broker' = broker \cup {m}
    /\ sent' = sent \cup {m.id}
    /\ UNCHANGED <<db, outbox, acked, allMsgs, nextMsgId>>

AckFromBroker(m) ==
    /\ m \in broker
    /\ acked' = acked \cup {m.id}
    /\ UNCHANGED <<db, outbox, broker, sent, allMsgs, nextMsgId>>

CleanupOutbox(m) ==
    /\ m \in outbox
    /\ m.id \in acked
    /\ outbox' = outbox \ {m}
    /\ UNCHANGED <<db, broker, sent, acked, allMsgs, nextMsgId>>
\end{lstlisting}
\end{primer}